
本章介绍 $J$-全纯球面在辛拓扑中的应用。我们将详细证明若干基本结果，并指出进一步
发展的方向。这里讨论的许多应用都源于
\personlater{geluomofu}{格罗莫夫}{Gromov}\srcite{160}，不过我们并不总是采用他的证明。
这个领域近年来发展极为活跃，因而不可能在此面面俱到。

第一节讨论哈密顿系统的周期轨道。我们证明：\BookRecordIndex{T0058}半正\BookRecordIndex{T0358}辛流形上的每个哈密顿系统都存在一条
周期轨道。在\BookRecordIndex{T0355}辛非球面的情形，我们重新得到 Schwarz 的一个定理：除非时间--$1$ 映射
是恒等映射，否则存在两条作用量不同的周期轨道；并说明如何用这一结果证明 \BookRecordIndex{T0027}Hofer 度量
的非退化性。第~\ref{sec:9-2} 节讨论\BookRecordIndex{T0219}拉格朗日子流形。我们证明
\personlater{geluomofu}{格罗莫夫}{Gromov}关于两件事的定理：一是存在边界落在\BookRecordIndex{T0219}拉格朗日
子流形上的非常值全纯圆盘，二是\BookRecordIndex{T0219}拉格朗日子流形与其任意哈密顿形变的交非空。所有这些
结果的证明都基于同一个想法：从一族\BookRecordIndex{T0209}柯西--黎曼方程形变之解的极限中，构造出一个有意义
的几何对象，例如周期轨道、$J$-全纯圆盘、\BookRecordIndex{T0131}弗洛尔连接轨道或拉格朗日交点。这里的形变
就是第~8 章所讨论的\BookRecordIndex{T0157}哈密顿扰动；不过相应的丛是平凡的，而且我们力求使不熟悉该章全部
细节的读者也能读懂这里的讨论。

接着，我们给出\BookRecordIndex{T0120}非挤压定理的一种证明，它使用吹起方法而不是原来的\BookRecordIndex{T0175}极小曲面论证。
第~\ref{sec:9-4} 节依循 \personlater{geluomofu}{格罗莫夫}{Gromov}\srcite{160} 与 McDuff
\srcite{258} 讨论有理曲面和直纹曲面的结构，第~\ref{sec:9-5} 节则证明
\personlater{geluomofu}{格罗莫夫}{Gromov}关于
$\mathbb R^4$、$\mathbb CP^2$ 和 $S^2\times S^2$ 的辛同胚群之结果。第~\ref{sec:9-6} 节介绍
\BookRecordIndex{T0003}$J$-全纯曲线技巧在 Hofer 几何中的若干应用。这一方向最早的工作是 Lalonde--McDuff
\srcite{223,224}。这里我们解释 Polterovich \srcite{325,326,327} 与 Seidel
\srcite{363,364} 关于 Seidel 同态的一项应用。最后，第~\ref{sec:9-7} 节讨论 Ruan \srcite{340}
和 McDuff \srcite{254} 的一些例子，其中用 $J$-全纯球面来区分\BookRecordIndex{T0358}辛流形。

这里讨论的应用限于亏格零情形，并且在精神上仍很接近
\personlater{geluomofu}{格罗莫夫}{Gromov}最初的想法。我们主要关注不必计数全纯曲线
数目（或者更准确地说，不必计数到超过一）的情形，而是从一张全纯圆盘或球面的存在性、
或者从这样一族曲线的存在性中提取信息，其中每一点都通过一条曲线。在这些情形中，无须
使用第~10 章的\BookRecordIndex{T0256}粘合定理。相反，对全纯球面计数所得的信息可以组织成\BookRecordIndex{T0229}量子上同调
（第~11 章），而这确实需要\BookRecordIndex{T0256}粘合定理。Floer 研究了圆柱、条带之类的非紧
\BookRecordIndex{T0003}$J$-全纯曲线，由此产生了如今所谓的\BookRecordIndex{T0132}弗洛尔同调，并带来许多卓越的应用。（概览见
第~12 章。）辛弗洛尔理论的出现早于\BookRecordIndex{T0229}量子上同调的发展，但二者联系紧密；它还需要另一类
\BookRecordIndex{T0256}粘合定理，而本书并不讨论这种定理。在四维情形及其接触几何中的三维对应物里，近年来
出现了基于非紧 \BookRecordIndex{T0003}$J$-全纯曲线研究的许多有趣进展，这些曲线可以拼合成有限能量叶状结构。
这一方法由 Hofer \srcite{178} 和 Hofer--Wysocki--Zehnder \srcite{183} 开创，Hind
\srcite{174} 与 Wendl \srcite{418} 又给出了更多有趣应用。本书不涉及这些结果，但它们
在精神上与本章介绍的若干思想十分接近。

\section{哈密顿系统的周期轨道}\label{sec:9-1}
\zhindex{hamidunxitongdezhouqiguidao}{哈密顿系统的周期轨道}

哈密顿系统周期解的研究有着悠久而精彩的历史，曾吸引 Weierstrass、Poincaré 和 Moser
等数学家。尤其是存在性问题，推动了辛拓扑中的许多进展；这一课题的文献浩如烟海，我们
不可能在这里展开讨论。我们只关注如下存在性定理：若不存在 Chern 数为负的
$J$-全纯球面，则周期哈密顿系统有 $1$-周期解。这个论证只给出一条周期轨道（附加一个
假设时给出两条），应把它看作一个原型，用来展示椭圆方法为何与这一问题相关。我们的
证明研究\BookRecordIndex{T0209}柯西--黎曼方程越来越大的扰动，因而与
\personlater{geluomofu}{格罗莫夫}{Gromov}\srcite{160} 的一些论证相似，但几何背景
不同：\personlater{geluomofu}{格罗莫夫}{Gromov}考虑边界落在\BookRecordIndex{T0219}拉格朗日子流形上的圆盘，
这里则考虑受扰 $J$-全纯球面，它们在极限中产生圆柱上的 弗洛尔方程之解。此外，
\personlater{geluomofu}{格罗莫夫}{Gromov}的结论稍弱：他只断言哈密顿流的时间--$1$
映射存在不动点，而不是存在可缩的 $1$-周期轨道。尽管如此，他的论文仍是这里所述思想的
基础。论文发表数年后，Floer \srcite{113,114,116} 发展出一种新的同调理论，即\BookRecordIndex{T0132}弗洛尔
同调，从而提供了计数周期轨道的方法。

设 $(M,\omega)$ 是紧\BookRecordIndex{T0358}辛流形，且
\[
  \mathbb R\times M\longrightarrow\mathbb R,
  \qquad (t,x)\longmapsto H_t(x)
\]
是光滑哈密顿函数，对每个 $t$ 都有 $H_t=H_{t+1}$。\sourcefootnote{原书把该哈密顿函数的类型误印为 $\mathbb R\to M$；但输入是
$(t,x)$，且 $H_t(x)$ 为实值，故其源、靶唯一应为
$\mathbb R\times M\to\mathbb R$。}
以 $X_{H_t}\in\operatorname{Vect}(M)$ 表示 $H_t$ 对应的\BookRecordIndex{T0160}哈密顿向量场，即
\[
  \iota(X_{H_t})\omega=dH_t,
\]
并考虑微分方程
\begin{equation}\label{eq:9-1-1}
  \dot x(t)=X_{H_t}(x(t)).
\end{equation}
若式~\eqref{eq:9-1-1} 的解 $x:\mathbb R\to M$ 对每个 $t$ 都满足
$x(t)=x(t+1)$，便称它是 $1$-周期的。用 $\mathcal P(H)$ 表示式~\eqref{eq:9-1-1}
的 $1$-周期解之集，用 $\mathcal P_0(H)$ 表示其中可缩的 $1$-周期解之集。

\begin{theorem}\label{thm:9-1-1}
设 $(M,\omega)$ 是闭连通\BookRecordIndex{T0358}辛流形。假设存在一个被 $\omega$ \BookRecordIndex{T0372}驯服的\BookRecordIndex{T0088}殆复结构
$J\in\mathcal J_\tau(M,\omega)$，使每个 $J$-全纯球面都有非负 Chern 数。则对于每个
哈密顿函数 $H:\mathbb R/\mathbb Z\times M\to\mathbb R$，方程
\eqref{eq:9-1-1} 都有一个可缩的 $1$-周期解。
\end{theorem}

每个哈密顿辛同胚 $\phi$ 都由某个 $1$-周期哈密顿函数生成，所以该定理蕴含：哈密顿群
$\operatorname{Ham}(M,\omega)$ 的每个元素 $\phi$ 都有不动点 $x$。（这就是
\personlater{geluomofu}{格罗莫夫}{Gromov}在 \srcite{160} 中对 $\omega$ 在
$\pi_2(M)$ 上为零的流形所证明的结果。）进一步，若
$\{\phi_t\}_{0\le t\le1}$ 是 $\operatorname{Ham}(M,\omega)$ 中从恒等映射到 $\phi$
的任意道路，则可选取 $x$，使环路 $t\mapsto\phi_t(x)$ 在 $M$ 中可缩。这带来如下有些
出人意料的推论。

\begin{corollary}\label{cor:9-1-2}
设 $(M,\omega)$ 满足定理~\ref{thm:9-1-1} 的假设，并设
\[
  \mathbb R/\mathbb Z\longrightarrow\operatorname{Ham}(M,\omega),
  \qquad t\longmapsto\phi_t
\]
是哈密顿辛同胚环路。则对每个 $x\in M$，环路
$\mathbb R/\mathbb Z\to M:t\mapsto\phi_t(x)$ 都可缩。换言之，由
$\{\phi_t\}\mapsto\{\phi_t(x)\}$ 给出的\BookRecordIndex{T0283}求值映射
\[
  \pi_1(\operatorname{Ham}(M,\omega))\longrightarrow\pi_1(M,x)
\]
是平凡的。
\end{corollary}

\begin{proof}
由定理~\ref{thm:9-1-1}，环路
$\mathbb R/\mathbb Z\to M:t\mapsto\phi_t(x)$ 中至少有一个可缩。因此它们全都可缩。
\end{proof}

这个推论意味着，对这样的流形 $(M,\omega)$，可以给哈密顿辛同胚 $\phi$ 的每个不动点
$x_0$ 指派 $\pi_1(M,x)$ 中的一个共轭类：它就是环路
$\{\phi_t(x_0)\}_{t\in[0,1]}$ 的自由同伦类，其中 $\phi_t$ 是时间--$1$ 映射为 $\phi$
的任意哈密顿道路。特别地，每个这样的 $\phi$ 都有一个良好定义的可缩不动点集
$\operatorname{Fix}_0(\phi)$，它由使上述环路可缩的点组成。

这里是用解析方法证明了一个纯拓扑结果。尽管这一推论另有证明（例如见 \srcite{228}），
但据我们所知，所有这些证明都以某种形式使用了 \BookRecordIndex{T0003}$J$-全纯曲线理论。

\paragraph{Arnold 猜想。}
在证明定理~\ref{thm:9-1-1} 之前，先把这个结果置于更大的背景中。Arnold
\srcite{16} 在 20 世纪 60 年代猜想：对每个哈密顿函数
$H:\mathbb R/\mathbb Z\times M\to\mathbb R$，都有
\begin{equation}\label{eq:9-1-2}
  \#\mathcal P_0(H)\ge \operatorname{Crit}(M),
\end{equation}
其中 $\operatorname{Crit}(M)$ 表示函数 $f:M\to\mathbb R$ 的\BookRecordIndex{T0233}临界点数目之最小值。这里
既可以允许任意的 $H$ 和 $f$，也可以限制在二者都非退化的情形。（这意味着 $f$ 是
Morse 函数，而且 $H$ 的周期轨道附近线性化流的时间--$1$ 映射没有等于 $1$ 的特征值。）
这一形式的猜想至今仍未解决。

这个猜想的同调版本更容易处理。在非退化情形，它断言：只要所有可缩的 $1$-周期解都
非退化，就有
\begin{equation}\label{eq:9-1-3}
  \#\mathcal P_0(H)\ge\dim H^*(M;\mathbb Q).
\end{equation}
这个结果最早由 Eliashberg \srcite{98} 对黎曼曲面证明，并由 Conley--Zehnder
\srcite{75} 对标准 $2n$-维环面证明。Floer \srcite{112,113,114,116} 的工作带来了
突破：他对单调\BookRecordIndex{T0358}辛流形证明了非退化 Arnold 猜想。此后关于这个问题的所有工作都建立在
Floer 的思想之上。Hofer--Salamon \srcite{180} 与 Ono \srcite{313} 把结果推广到\BookRecordIndex{T0058}半正
情形；第 12.1 节将概述他们的证明。现在 Fukaya--Ono \srcite{127} 与 Liu--Tian
\srcite{249} 已对所有紧\BookRecordIndex{T0358}辛流形证明了式~\eqref{eq:9-1-3} 形式的非退化 Arnold 猜想。
因此，定理~\ref{thm:9-1-1} 中关于全纯球面 Chern 数的假设可以去掉。

\begin{remark}\label{rem:9-1-3}
\begin{enumerate}[label=\textnormal{(\roman*)}]
\item 当哈密顿函数 $H$ 与 $t$ 无关时，每个 $H$ 的\BookRecordIndex{T0233}临界点都给出
$\mathcal P_0(H)$ 的一个元素，故估计~\eqref{eq:9-1-2} 可由 Morse 理论\BookRecordIndex{T0326}推出。当 $M$
的奇数维同调为零时，它也可由 Lefschetz 不动点公式\BookRecordIndex{T0326}推出。进一步，若 $M$ 的 Euler
示性数非零，则 Lefschetz 不动点公式蕴含 $H$ 的流之时间--$1$ 映射有不动点；不过它并不
断言这个不动点是可缩的，也就是说，它未必对应可缩环路。

\item （Conley 猜想）Conley 猜想 $2n$-维环面上的每个哈密顿系统都有无穷多条两两不同、
周期为整数的周期解。这样的断言应当对更一般的流形成立，但不是对所有流形都成立；
例如 $S^2$ 上的无理旋转只有两个周期点。Salamon--Zehnder \srcite{356} 在假设
$[\omega]$ 在 $\pi_2(M)$ 上为零、且每个 $1$-周期解至少有一个不等于 $1$ 的 Floquet
乘子时，证明了无穷多条周期轨道的存在性。完整的猜想如今已经解决：Hingston
\srcite{176} 首先解决了环面情形，Ginzburg \srcite{143} 随后解决了任意闭\BookRecordIndex{T0355}辛非球面
情形。不过，它成立的最大流形类别仍不清楚。

\item （Weinstein 猜想）这个猜想最初的形式断言：辛 Euclidean 空间中任意闭\BookRecordIndex{T0187}接触型超
曲面上的特征流都有周期轨道。Viterbo \srcite{402} 在一篇重要的早期论文中证实了这个
猜想；该工作在 Conley--Zehnder 关于环面 Arnold 猜想的工作 \srcite{75} 与
Ekeland--Hofer 的辛容量变分理论 \srcite{97} 之间建立了关键联系。下面给出的
定理~\ref{thm:9-1-1} 之证明所用的思想，后来由 Hofer--Viterbo \srcite{182}、
Liu--Tian \srcite{251} 和 Lu \srcite{253} 推广，从而在相当一般的假设下证明了周期
轨道的存在性。Frauenfelder 与 Schlenk 的精致论文 \srcite{123} 以及 Albers--Hofer
\srcite{12} 中可以找到这个课题的一些新结果。Taubes \srcite{394} 用后来又被他推广
以构造 Seiberg--Witten--Floer 理论与 Hutchings--Sullivan \BookRecordIndex{T0208}嵌入接触同调之间同构的论证，
成功对三维接触流形上的 Reeb 流证明了该猜想。不过，并非 Euclidean 空间中的每个闭超
曲面都容许这样的轨道；Ginzburg \srcite{141,142} 与 Herman \srcite{172,424} 找到了
有趣的反例。
\end{enumerate}
\end{remark}

定理~\ref{thm:9-1-1} 的证明建立在二维球面上的非线性\BookRecordIndex{T0209}柯西--黎曼方程及第~\ref{sec:8-1} 节所讨论
类型的扰动之上。我们假设读者熟悉该节。受扰方程形如
\begin{equation}\label{eq:9-1-4}
  \bar\partial_J(v)+X_K(v)^{0,1}=0,
\end{equation}
其中 $J\in C^\infty(S^2,\mathcal J_\tau(M,\omega))$，且
$K\in\Omega^1(S^2,C_0^\infty(M))$。这里 $C_0^\infty(M)$ 表示均值为零的光滑函数空间。
由引理~\ref{lem:8-1-6} 和注~\ref{rem:8-1-7}，式~\eqref{eq:9-1-4} 的解
$v:S^2\to M$ 的能量满足
\begin{equation}\label{eq:9-1-5}
  E_K(v):=\int_{S^2}|dv+X_K(v)|_J^2\,\mathrm{dvol}_{S^2}
  \le \int_{S^2}v^*\omega+\|R_K\|,
\end{equation}
其中 $\|R_K\|$ 表示曲率的 \BookRecordIndex{T0028}Hofer 范数。回忆
$\mathcal J_+(M,\omega;\kappa)$ 的定义：它由所有这样的
$\omega$-\BookRecordIndex{T0372}驯服\BookRecordIndex{T0088}殆复结构 $J\in\mathcal J_\tau(M,\omega)$ 组成，即每个满足
$E(v)\le\kappa$ 的 $J$-全纯球面 $v:S^2\to M$ 都有非负 Chern 数。由
引理~\ref{lem:6-4-7}，这是 $\mathcal J_\tau(M,\omega)$ 的开子集（可能为空）。

\begin{proposition}\label{prop:9-1-4}
设 $(M,\omega)$ 是紧\BookRecordIndex{T0358}辛流形，$J=\{J_z\}_{z\in S^2}$ 是一族光滑的
$\omega$-\BookRecordIndex{T0372}驯服\BookRecordIndex{T0088}殆复结构，且 $K\in\Omega^1(S^2,C_0^\infty(M))$。若存在实数
$\kappa>\|R_K\|$，使每个 $z$ 都有 $J_z\in\mathcal J_+(M,\omega;\kappa)$，则
式~\eqref{eq:9-1-4} 存在一个解 $v:S^2\to M$，其能量满足
\[
  E_K(v)\le\|R_K\|.
\]
\end{proposition}

在证明命题之前，我们先用它证明主要结果。

\begin{proof}[定理~\ref{thm:9-1-1} 的证明]
设 $J$ 是 $M$ 上被 $\omega$ \BookRecordIndex{T0372}驯服的\BookRecordIndex{T0088}殆复结构，并且每个 $J$-全纯球面都有非负 Chern 数；
这样的\BookRecordIndex{T0088}殆复结构由假设存在。对每个 $T>0$，取\BookRecordIndex{T0188}截断函数
$\beta_T:\mathbb R\to[0,1]$，使它在 $[-T,T]$ 上等于一、支集包含在
$[-T-1,T+1]$ 中，并且在 $[T,T+1]$ 上非增、在 $[-T-1,-T]$ 上非减。给定光滑
哈密顿函数
\[
  \mathbb R/\mathbb Z\longrightarrow C_0^\infty(M),\qquad t\longmapsto H_t,
\]
考虑受扰\BookRecordIndex{T0209}柯西--黎曼方程
\begin{equation}\label{eq:9-1-6}
  \partial_su+J(u)\bigl(\partial_tu-\beta_T(s)X_{H_t}(u)\bigr)=0
\end{equation}
及其光滑映射解 $u:\mathbb R\times\mathbb R/\mathbb Z\to M$。通过球极投影把
$S^2$ 与 Riemann 球面 $\mathbb C\cup\{\infty\}$ 等同，并通过全纯微分同胚
\[
  \phi:\mathbb R\times\mathbb R/\mathbb Z\longrightarrow
  \mathbb C\setminus\{0\}\subset S^2,
  \qquad \phi(s,t):=e^{2\pi(s+it)}
\]
把圆柱与 $S^2\setminus\{0,\infty\}$ 等同。由\BookRecordIndex{T0214}可去奇点定理
（定理~\ref{thm:4-1-2}），式~\eqref{eq:9-1-6} 的解 $u$ 具有有限能量
\[
  E(u):=\int_0^1\int_{-\infty}^{\infty}|\partial_su|^2\,ds\,dt<\infty
\]
当且仅当函数 $v_u:=u\circ\phi^{-1}:S^2\setminus\{0,\infty\}\to M$ 能延拓为
$S^2\to M$ 的光滑映射。此外，$v_u$ 是式~\eqref{eq:9-1-4} 的解，其中扰动
$K=K_T$ 在圆柱上的拉回为
$\phi^*K_T=-\beta_T(s)H_t\,dt$。关键观察是，$K_T$ 曲率的 \BookRecordIndex{T0028}Hofer 范数为
\[
  \|R_{K_T}\|
  =\int_{-\infty}^{\infty}\int_0^1|\dot\beta_T(s)|
    \left(\max_MH_t-\min_MH_t\right)\,dt\,ds
  =2\|H\|.
\]
因此，由命题~\ref{prop:9-1-4}，对每个 $T>0$ 都存在式~\eqref{eq:9-1-6}
的一个解 $u_T:\mathbb R\times\mathbb R/\mathbb Z\to M$，并且
$E(u_T)\le2\|H\|$。把 $u_T$ 限制在子集
$[-T,T]\times\mathbb R/\mathbb Z$ 上，\sourcefootnote{原书此处把圆柱的第二因子误印为
$\mathbb R$；但 $u_T$ 的定义域及随后对 $t\in[0,1]$ 的积分都唯一确定该因子应为
$\mathbb R/\mathbb Z$。}可得
\[
  \int_0^1\int_{-T}^{T}
  |\partial_tu_T-X_{H_t}(u_T)|^2\,ds\,dt
  \le E(u_T)\le2\|H\|.
\]
所以存在 $s_T\in[-T,T]$，使环路 $x_T(t):=u_T(s_T,t)$ 满足
\[
  \int_0^1|\dot x_T(t)-X_{H_t}(x_T(t))|^2\,dt\le\frac{\|H\|}{T}.
\]
由于 $M$ 紧，存在序列 $T_\nu\to\infty$，使 $x_{T_\nu}$ 在 $W^{1,2}$ 拓扑下收敛到
式~\eqref{eq:9-1-1} 的一个 $1$-周期解
$x:\mathbb R/\mathbb Z\to M$。每个 $x_{T_\nu}$ 都可缩，因为它是定义域为二维球面的映射
$v_{T_\nu}$ 的限制。因此 $x$ 也可缩。定理得证。
\end{proof}

\begin{proof}[命题~\ref{prop:9-1-4} 的证明]
我们先证明一个稍弱、但已足够用于这里应用的结果。固定常数 $\kappa>0$，使命题的假设
成立。于是 $\mathcal J_+(M,\omega;\kappa)$ 非空，并由引理~\ref{lem:6-4-7} 知它是
$\mathcal J_\tau(M,\omega)$ 的开子集。固定
$J_0\in\mathcal J_+(M,\omega;\kappa)$，并以 $\mathcal J_0(\kappa)$ 表示所有光滑映射
\[
  S^2\longrightarrow\mathcal J_+(M,\omega;\kappa),\qquad z\longmapsto J_z
\]
所成的空间，其中这些映射与常值映射 $J_z\equiv J_0$ 同伦。以
$\mathcal H_0(\kappa)$ 表示所有满足如下条件的\BookRecordIndex{T0157}哈密顿扰动
$K\in\Omega^1(S^2,C_0^\infty(M))$ 之集：在微分同胚
\[
  \phi:\mathbb R\times\mathbb R/\mathbb Z\longrightarrow S^2\setminus\{0,\infty\},
  \qquad \phi(s,t):=e^{2\pi(s+it)}
\]
\sourcefootnote{原书把 $\phi$ 称为到整个 $S^2$ 的微分同胚；但该指数映射的像是
$S^2\setminus\{0,\infty\}$，且前文也正是用它等同圆柱与双穿孔球面，故补正其靶。}
下，$K$ 的拉回形如 $\phi^*K=G\,dt$，且曲率的 \BookRecordIndex{T0028}Hofer 范数满足
\[
  \|R_K\|=
  \int_{\mathbb R\times\mathbb R/\mathbb Z}
  \left(\max_M\partial_sG-\min_M\partial_sG\right)\,ds\,dt<\kappa.
\]
注意 $\mathcal H_0(\kappa)$ 可缩：若 $K\in\mathcal H_0(\kappa)$，则对
$0\le\lambda\le1$ 有 $\lambda K\in\mathcal H_0(\kappa)$。

与 $J\in\mathcal J_0(\kappa)$ 和 $K\in\mathcal H_0(\kappa)$ 相联系的是
式~\eqref{eq:9-1-4} 的可缩解 $v:S^2\to M$ 的\BookRecordIndex{T0246}模空间
\[
  \mathcal M_0(J,K):=
  \bigl\{v:S^2\to M\bigm|v\text{ 满足式~\eqref{eq:9-1-4}，且 }v\text{ 可缩}\bigr\}.
\]
我们断言，对 $\mathcal J_0(\kappa)\times\mathcal H_0(\kappa)$ 中稠密的偶对
$(J,K)$，$\mathcal M_0(J,K)$ 非空。证明依赖\BookRecordIndex{T0249}能量恒等式~\eqref{eq:9-1-5}；在当前
情形，它化为
\[
  E_K(v)\le\|R_K\|<\kappa.
\]
以 $\mathcal J\mathcal H_{\mathrm{reg}}(\kappa)$ 表示所有偶对
$(J,K)\in\mathcal J_0(\kappa)\times\mathcal H_0(\kappa)$ 中，在定义
~\ref{def:8-5-2} 的意义下对 $k=\#I=1$ 和点 $w=0$ 正则者之集。定理
~\ref{thm:8-5-3} 断言，该集合在
$\mathcal J_0(\kappa)\times\mathcal H_0(\kappa)$ 中是剩余的。若
$(J,K)\in\mathcal J\mathcal H_{\mathrm{reg}}(\kappa)$，则由定理~\ref{thm:8-5-1}，
\BookRecordIndex{T0283}求值映射
\[
  \operatorname{ev}_{J,K}:\mathcal M_0(J,K)\longrightarrow M,
  \qquad \operatorname{ev}_{J,K}(u):=u(0)
\]
是维数 $2n=\dim M$ 的\BookRecordIndex{T0337}伪循环。由于 $\mathcal J_0(\kappa)$ 与
$\mathcal H_0(\kappa)$ 都连通，它所表示的 $H_{2n}(M)$ 中同调类与 $J,K$ 无关。相应的
\BookRecordIndex{T0143}格罗莫夫--威滕不变量为
\[
  \operatorname{GW}_{0,1}^M(a)
  =\operatorname{ev}_{J,K}\mathbin{\boldsymbol\cdot}\mathrm{pt}=1,
  \qquad a:=\operatorname{PD}([\mathrm{pt}])\in H^{2n}(M).
\]
这里把\BookRecordIndex{T0283}求值映射写成 $\operatorname{ev}_{J,K}$。
\sourcefootnote{原书在本式及下文一般偶对的对应公式中两次误印为
$\operatorname{ev}_{J,H}$；本段扰动变量始终是 $K$，且紧邻处定义的也是
$\operatorname{ev}_{J,K}$，故据此补正。}
为看出不变量等于一，考虑零哈密顿量。由引理~\ref{lem:6-7-6}，每个 $K=0$ 的偶对
$(J,K)$ 都属于正则集 $\mathcal J\mathcal H_{\mathrm{reg}}(\kappa)$；而对每个
$x_0\in M$，式~\eqref{eq:9-1-4} 在 $K=0$ 时恰有一个满足 $v(0)=x_0$ 的可缩解，即
常值解 $v(z)\equiv x_0$。所以对每个
$(J,K)\in\mathcal J\mathcal H_{\mathrm{reg}}(\kappa)$，都有
$\mathcal M_0(J,K)\ne\varnothing$。

现在取 $J\in\mathcal J_0(\kappa)$、$K\in\mathcal H_0(\kappa)$，并选取在
$C^\infty$ 拓扑下收敛到 $(J,K)$ 的正则序列
$(J_\nu,K_\nu)\in\mathcal J\mathcal H_{\mathrm{reg}}(\kappa)$。由刚才所证，
对每个 $\nu$ 都有 $\mathcal M_0(J_\nu,K_\nu)\ne\varnothing$；任取
$v_\nu\in\mathcal M_0(J_\nu,K_\nu)$。由于 $K_\nu$ 的曲率一致有界，由
注~\ref{rem:8-1-7}，乘积流形 $S^2\times M$ 上存在一个\BookRecordIndex{T0363}辛形式
$\omega_{K,\kappa}$，对每个 $\nu$ 都\BookRecordIndex{T0372}驯服由 $J_\nu,K_\nu$ 诱导的\BookRecordIndex{T0088}殆复结构
$\widetilde J_\nu$（见习题~\ref{ex:8-1-4}）。因此，可以把定理~\ref{thm:4-6-1}
用于 $S^2\times M$ 中 $v_\nu$ 的图像序列，得到有限集 $Z\subset S^2$ 及一个子列；
该子列在 $S^2\setminus Z$ 的每个紧子集上连同所有阶导数一致收敛到式
~\eqref{eq:9-1-4} 的解 $v:S^2\to M$。因为
$E_{K_\nu}(v_\nu)\le\|R_{K_\nu}\|$ 且 $R_{K_\nu}\to R_K$，极限解满足所需的能量
界 $E_K(v)\le\|R_K\|$。

至此，我们对
$\mathcal J_0(\kappa)\times\mathcal H_0(\kappa)$ 中的偶对 $(J,K)$ 证明了命题；
这已经足以证明定理~\ref{thm:9-1-1}。对一般偶对，可以如下推广上述论证。固定
$\kappa>0$，用 $\mathcal H(\kappa)$ 表示满足 $\|R_K\|<\kappa$ 的所有\BookRecordIndex{T0157}哈密顿扰动
$K\in\Omega^1(S^2,C_0^\infty(M))$ 之集，用 $\mathcal J(\kappa)$ 表示所有光滑映射
\[
  S^2\longrightarrow\mathcal J_+(M,\omega;\kappa),\qquad z\longmapsto J_z
\]
所成的空间。仍以 $\mathcal J\mathcal H_{\mathrm{reg}}(\kappa)$ 表示所有满足定义
~\ref{def:8-5-2} 中 $k=\#I=1$、$w=0$ 之条件的偶对
$(J,K)\in\mathcal J(\kappa)\times\mathcal H(\kappa)$。这一次，我们不知道
$\mathcal J(\kappa)\times\mathcal H(\kappa)$ 是否连通，因此在它的每个连通分支上得到
一个\BookRecordIndex{T0143}格罗莫夫--威滕不变量
\[
  \operatorname{GW}_{0,1}^M(\operatorname{PD}([\mathrm{pt}]))
  =\operatorname{ev}_{J,K}\mathbin{\boldsymbol\cdot}\mathrm{pt}.
\]
我们断言，该不变量在每个连通分支上都等于一。前面的论证适用于每个包含形如 $(J,0)$
之偶对的连通分支。现在设 $R_K=0$。于是存在哈密顿规范变换
\[
  S^2\longrightarrow\operatorname{Ham}(M,\omega),\qquad z\longmapsto\psi_z,
\]
使 $\psi^*K=0$。因此，对每个包含某个满足 $R_K=0$ 的偶对 $(J,K)$ 的连通分支，
上述不变量都等于一。由下面的习题~\ref{ex:9-1-5}，
$\mathcal J(\kappa)\times\mathcal H(\kappa)$ 的每个连通分支都包含这样的偶对，故不变量
总等于一。于是，对每个正则偶对
$(J,K)\in\mathcal J\mathcal H_{\mathrm{reg}}(\kappa)$，都有
$\mathcal M_0(J,K)\ne\varnothing$。最后采用与上面相同的逼近论证，便证明了命题
~\ref{prop:9-1-4}。
\end{proof}

证明的后一部分稍显笨重，原因是本书并未对所有紧\BookRecordIndex{T0358}辛流形定义\BookRecordIndex{T0143}格罗莫夫--威滕不变量。
因此，我们只能使用不存在负 Chern 数 $J$-全纯球面的\BookRecordIndex{T0088}殆复结构 $J$。所得空间
$\mathcal J(\kappa)$ 未必连通，于是还需要额外论证，证明每个连通分支上的不变量都等于
一。若采用\BookRecordIndex{T0143}格罗莫夫--威滕不变量的一般定义，这些额外论证就不再需要，命题
~\ref{prop:9-1-4} 中的假设 $J_z\in\mathcal J_+(M,\omega;\kappa)$ 也可以去掉。

\begin{exercise}\label{ex:9-1-5}
设 $K\in\Omega^1(S^2,C_0^\infty(M))$。证明存在光滑道路
\[
  [0,1]\longrightarrow\Omega^1(S^2,C_0^\infty(M)),\qquad
  \lambda\longmapsto K_\lambda,
\]
使得
\[
  R_{K_0}=0,\qquad K_1=K,\qquad
  \sup_{0\le\lambda\le1}\|R_{K_\lambda}\|\le\|R_K\|.
\]

\emph{提示：}考虑在嵌入
$\phi:\mathbb R\times\mathbb R/\mathbb Z\to S^2$ 下 $K$ 到
$\mathbb R\times\mathbb R/\mathbb Z$ 的拉回
\[
  \phi^*K=F\,ds+G\,dt.
\]
若 $F$ 或 $G$ 为零，结论很容易：只须考虑族 $\lambda K$。用一族规范变换同伦 $K$，使
$s\le1$ 时 $F_{s,t}=0$；再作同伦，使 $s\le0$ 时 $F=G=0$。然后对 $s\ge0$ 重复同样
的论证。注意，这个结果成立只是因为底\BookRecordIndex{T0072}层的\BookRecordIndex{T0156}哈密顿纤维化是平凡的。第~\ref{sec:9-6} 节将说明，
把 $\|R\|$ 做到任意小存在拓扑障碍。
\end{exercise}

\paragraph{辛作用量。}
下面讨论这些思想的一些推广，其中会用到辛作用泛函。最重要的结果只适用于如下情形：
$(M,\omega)$ 是紧连通\BookRecordIndex{T0358}辛流形，而且 $\omega$ 的上同调类在 $\pi_2(M)$ 上为零，即对每个
光滑映射 $v:S^2\to M$ 都有
\[
  \int_{S^2}v^*\omega=0.
\]
这样的流形称为\emph{\BookRecordIndex{T0355}辛非球面}。这个假设具有真正的几何意义，因为没有它时，作用泛函
只能定义在环路空间的一个覆叠上。此外，它蕴含不存在 $J$-全纯球面，因而证明在技术上
更简单。

以 $L_0M$ 表示可缩环路 $x:\mathbb R/\mathbb Z\to M$ 的空间，并以
$\widetilde{L_0M}$ 表示如下覆叠：其元素是偶对 $(x,[u])$，其中 $x\in L_0M$，而
$[u]$ 是 $x$ 到圆盘上的一个延拓。更准确地说，$[u]$ 表示光滑映射
\[
  u:\mathbb D\longrightarrow M,
  \qquad \mathbb D:=\{z\in\mathbb C\mid |z|\le1\}
\]
相对于边界的同伦类，并要求对 $t\in\mathbb R$ 有
$u(e^{2\pi it})=x(t)$。光滑哈密顿函数
$H:\mathbb R/\mathbb Z\times M\to\mathbb R$ 决定辛作用泛函
\[
  \mathcal A_H:\widetilde{L_0M}\longrightarrow\mathbb R,
\]
其定义为
\begin{equation}\label{eq:9-1-7}
  \mathcal A_H(x,[u])
  :=-\int_{\mathbb D}u^*\omega-\int_0^1H_t(x(t))\,dt.
\end{equation}
这里的符号选取使得 $\mathcal A_H$ 与 Euclidean 空间中沿环路的经典作用积分
$p\,dq-H\,dt$ 一致。下面的引理~\ref{lem:9-1-8} 将说明，$\mathcal A_H$ 的\BookRecordIndex{T0233}临界点
正是这样的偶对 $(x,[u])$：其中 $x$ 是 $H$ 所生成的\BookRecordIndex{T0158}哈密顿同痕
$\{\phi_t^H\}_{0\le t\le1}$ 的 $1$-周期轨道，因而是式~\eqref{eq:9-1-1} 的解。
进一步，由引理~\ref{lem:9-1-9}，若把每个 $H_t$ 归一化为均值零，则\BookRecordIndex{T0234}临界值
$\mathcal A_H(x,[u])$ 只依赖于哈密顿道路
$\{\phi_t^H\}_{0\le t\le1}$ 相对于端点的同伦类。（为了使这句话有意义，$u$ 必须随
$H$ 连续变化。）所以\BookRecordIndex{T0234}临界值集
\[
  \operatorname{Spec}(H)
  :=\{\mathcal A_H(x,[u])\mid x\in\mathcal P_0(H)\}
\]
只依赖于 $H$ 在哈密顿群的万有覆叠
$\widetilde{\operatorname{Ham}}(M,\omega)$ 中所定义的元素 $\widetilde\phi$。这里
$\widetilde\phi:=(\phi,[\phi_t])$，其中 $[\phi_t]$ 表示 $H$ 所生成的哈密顿流
$\{\phi_t\}_{0\le t\le1}$ 相对于固定端点的同伦类，而 $\phi$ 是时间--$1$ 映射
$\phi_1$。集合 $\operatorname{Spec}(H)$ 称为 $\widetilde\phi$ 的\emph{\BookRecordIndex{T0419}作用谱}
$\operatorname{Spec}(\widetilde\phi)$。Oh \srcite{305}、Schwarz \srcite{360} 和
Entov--Polterovich \srcite{105} 已经说明，它包含关于 $\widetilde\phi$ 性质的大量有趣
信息。进一步讨论见第 12.4 节。

若 $(M,\omega)$ \BookRecordIndex{T0355}辛非球面，则 $\mathcal A_H(x,[u])$ 与 $u$ 的选取无关，所以在这种
情形中将从记号里省略 $u$。对\BookRecordIndex{T0355}辛非球面流形，定理~\ref{thm:9-1-1} 可加强如下。

\begin{theorem}[Schwarz \srcite{360}]\label{thm:9-1-6}
设 $(M,\omega)$ 是紧连通\BookRecordIndex{T0355}辛非球面流形。设
$H:\mathbb R/\mathbb Z\times M\to\mathbb R$ 是光滑哈密顿函数，且它的时间--$1$
映射不等于恒等映射。则存在两条可缩周期轨道 $x,y\in\mathcal P_0(H)$，使得
\[
  -\int_0^1\max_MH_t\,dt
  \le\mathcal A_H(y)<\mathcal A_H(x)
  \le-\int_0^1\min_MH_t\,dt.
\]
\end{theorem}

Schwarz 在 \srcite{360} 中实际证明了：\BookRecordIndex{T0419}作用谱有两个特别指定的（并且互异的）元素
$\mathcal A_H(x_0)$ 和 $\mathcal A_H(y_0)$，它们的选取只依赖于 $H$ 的时间--$1$ 映射
$\phi$。这个更强的结果使用第 12.4 节构造的\BookRecordIndex{T0271}谱不变量，似乎无法由这里讨论的较初等方法
得到。他还证明，若 $H_t\in C_0^\infty(M)$，则作用量 $\mathcal A_H(x_0)$ 只依赖于
$x_0$ 和 $\phi$，而不依赖于从恒等映射到 $\phi$ 的\BookRecordIndex{T0158}哈密顿同痕。这同样是一个深刻的
结果。下面的习题~\ref{ex:9-1-11} 说明，两个不动点的相对作用量与\BookRecordIndex{T0158}哈密顿同痕无关。
这是一个初等观察，并且对任意流形都有类似结论。相比之下，绝对作用量的对应断言由推论
~\ref{cor:8-6-10} 得出，而后者又依赖 \BookRecordIndex{T0003}$J$-全纯曲线的\BookRecordIndex{T0256}粘合定理。此外，它只在非球面
情形成立。

Polterovich 在 \srcite{329} 中发现了定理~\ref{thm:9-1-6} 对\BookRecordIndex{T0358}辛流形上有限生成群作用
的卓越应用。例如，他证明了从 $\operatorname{SL}(3,\mathbb Z)$ 到正亏格黎曼曲面之
辛同胚群恒等连通分支的每个同态都必为平凡的。他的理论更一般地适用于任意维数的\BookRecordIndex{T0355}辛非
球面流形。

在开始证明定理~\ref{thm:9-1-6} 之前，我们再考察辛作用量的若干性质。先说明这里采用
的符号约定。

\begin{remark}[符号约定]\label{rem:9-1-7}
遗憾的是，常用文献中存在许多彼此不\BookRecordIndex{T0348}相容的符号约定。我们的选择受两条首要原则支配：
坐标为 $z=x+iy$ 的 $\mathbb R^2=\mathbb C$ 上，辛结构应为
$\omega=dx\wedge dy$；函数 $H:\mathbb R^2\to\mathbb R$ 所生成的哈密顿方程应为
\[
  \dot x=\frac{\partial H}{\partial y},
  \qquad
  \dot y=-\frac{\partial H}{\partial x}.
\]
因此，与这个流相切的向量场 $X_H$ 必须满足
\begin{equation}\label{eq:9-1-8}
  \omega(X_H,\,\cdot)=dH.
\end{equation}
（局部坐标中的哈密顿微分方程在文献里是标准的；不过许多作者同时改变
$\mathbb R^{2n}$ 上标准\BookRecordIndex{T0363}辛形式的符号以及式~\eqref{eq:9-1-8} 中的符号。）在上述约定
下，由 $H_t$ 生成的哈密顿辛同胚族 $\{\phi_t^H\}_{0\le t\le1}$ 满足
\begin{equation}\label{eq:9-1-9}
  \partial_t\phi_t^H(x)=X_{H_t}(\phi_t^H(x)).
\end{equation}
只要处于哈密顿动力学的范围内，我们都采用这些符号约定。不过，在辛\BookRecordIndex{T0279}纤维化上的\BookRecordIndex{T0155}哈密顿
联络这一几何背景中，为了使讨论与规范理论中的标准约定\BookRecordIndex{T0348}相容，我们选取的水平\BookRecordIndex{T0263}平行移动
方程是 $\partial_t\psi+X_{G_t}\circ\psi_t=0$，而不是式~\eqref{eq:9-1-9}。（见
第~\ref{sec:8-1}、\ref{sec:8-2} 和~\ref{sec:9-6} 节。）
\end{remark}

下面两个引理建立辛作用量最重要的性质。先说明 $\mathcal A_H$ 的\BookRecordIndex{T0233}临界点就是
式~\eqref{eq:9-1-1} 的可缩周期解。证明使用如下事实：由于 $\widetilde{L_0M}$ 是
$L_0M$ 的覆叠，它在 $(x,[u])$ 处的切空间可以与 $L_0M$ 在 $x$ 处的切空间等同。

\begin{lemma}\label{lem:9-1-8}
设 $(M,\omega)$ 是\BookRecordIndex{T0358}辛流形，并固定哈密顿函数
$H\in C^\infty(\mathbb R/\mathbb Z\times M)$。设
\[
  (x,[u])\in\widetilde{L_0M},
  \qquad
  \xi\in T_xL_0M=C^\infty(\mathbb R/\mathbb Z,x^*TM).
\]
则
\[
  d\mathcal A_H(x,[u])\xi
  =\int_0^1\omega\bigl(\dot x-X_{H_t}(x),\xi\bigr)\,dt.
\]
\end{lemma}

\begin{proof}
选取两个光滑映射
\[
  \mathbb R\times\mathbb R/\mathbb Z\longrightarrow M,
  \quad (s,t)\longmapsto x^s(t)=v(s,t),
  \qquad
  \mathbb R\times\mathbb D\longrightarrow M,
  \quad (s,z)\longmapsto u^s(z),
\]
使得
\[
  x^0(t)=x(t),\qquad u^0(z)=u(z),\qquad
  u^s(e^{2\pi it})=x^s(t),\qquad \partial_sv(0,t)=\xi(t).
\]
于是映射
$\mathbb R\to\widetilde{L_0M}:s\mapsto(x^s,[u^s])$ 是 $L_0M$ 中道路
$s\mapsto x^s$ 的提升，并通过 $(x^0,[u^0])=(x,[u])$。此外，
\[
  \mathcal A_H(x^{s_0},[u^{s_0}])
  =\mathcal A_0(x^0,[u^0])
  -\int_0^{s_0}\int_0^1\omega(\partial_sv,\partial_tv)\,dt\,ds
  -\int_0^1H_t(v(s_0,t))\,dt.
\]
所以
\[
  \frac d{ds}\mathcal A_H(x^s,[u^s])
  =-\int_0^1\omega\bigl(\partial_sv,\partial_tv-X_{H_t}(v)\bigr)\,dt,
\]
令 $s=0$ 即得结论。
\end{proof}

接着证明，\BookRecordIndex{T0419}作用谱 $\operatorname{Spec}(H)$ 只依赖于 $H$ 所生成的万有覆叠
$\widetilde{\operatorname{Ham}}$ 中的元素 $\widetilde\phi$。更准确地说，设
$(M,\omega)$ 是无边界紧连通\BookRecordIndex{T0358}辛流形，且
\[
  [0,1]\times\mathbb R/\mathbb Z\times M\longrightarrow\mathbb R,
  \qquad (\lambda,t,x)\longmapsto H_t^\lambda(x)
\]
是一族光滑的、归一化的 $1$-周期函数，因而对所有 $t,\lambda$ 都有
\begin{equation}\label{eq:9-1-10}
  \int_MH_t^\lambda\,\omega^n=0.
\end{equation}
对每个 $\lambda$，令 $\{\phi_t^\lambda\}_{t\in\mathbb R}$ 为 $H^\lambda$ 通过
\[
  \partial_t\phi_t^\lambda=X_t^\lambda\circ\phi_t^\lambda,
  \qquad \phi_0^\lambda=\operatorname{id},
  \qquad \iota(X_t^\lambda)\omega=dH_t^\lambda
\]
所生成的\BookRecordIndex{T0158}哈密顿同痕。假设时间--$1$ 映射 $\phi:=\phi_1^\lambda$ 与 $\lambda$ 无关。
下一个引理断言，在这些假设下，\BookRecordIndex{T0419}作用谱 $\operatorname{Spec}(H^\lambda)$ 与 $\lambda$
无关。

\begin{lemma}\label{lem:9-1-9}
设 $x_0\in\operatorname{Fix}_0(\phi)$ 是可缩不动点。对
$\lambda\in[0,1]$，定义环路
$x^\lambda:\mathbb R/\mathbb Z\to M$ 为
$x^\lambda(t):=\phi_t^\lambda(x_0)$，并设
\[
  [0,1]\times\mathbb D\longrightarrow M,
  \qquad (\lambda,z)\longmapsto u^\lambda(z)
\]
是任意光滑映射，满足对所有 $\lambda,t$ 都有
$u^\lambda(e^{2\pi it})=x^\lambda(t)$。若按式~\eqref{eq:9-1-10} 归一化
$H_t^\lambda$，则函数
\[
  \lambda\longmapsto\mathcal A_{H^\lambda}(x^\lambda,[u^\lambda])
\]
是常值函数。因此，每个
$\widetilde\phi\in\widetilde{\operatorname{Ham}}(M,\omega)$ 都有良好定义的\BookRecordIndex{T0419}作用谱。
\end{lemma}

\begin{proof}
对每个 $\lambda$，定义 $f^\lambda:M\to\mathbb R$ 为
\[
  f^\lambda(x):=-\int_{[0,1]^2}u_x^*\omega
  +\int_0^1\left(
    H_t^0(\phi_t^0(x))-H_t^\lambda(\phi_t^\lambda(x))
  \right)\,dt,
\]
其中 $u_x:[0,1]\times[0,1]\to M$ 定义为
$u_x(s,t):=\phi_t^{s\lambda}(x)$。注意，对每个 $x\in M$，$u_x$ 在方形边界上的限制
参数化环路 $\ell_x$：先沿道路 $\phi_t^\lambda(x)$ 从 $x$ 走到 $\phi(x)$，再沿
$\phi_{1-t}^0(x)$ 回到 $x$。这个环路由 $H_t^\lambda$ 与 $-H_{1-t}^0$ 拼接而成的
哈密顿函数 $K^\lambda$ 生成，故
\[
  f^\lambda(x)=\mathcal A_{K^\lambda}(\ell_x,[u_x]).
\]
因此，由引理~\ref{lem:9-1-8}，对所有 $x\in M$，$f^\lambda(x)$ 都是
$\mathcal A_{K^\lambda}:\widetilde{L_0M}\to\mathbb R$ 的\BookRecordIndex{T0234}临界值。由于 $M$ 连通，
$f^\lambda$ 必须是关于 $x$ 的常值函数。

现在对 $\lambda$ 求导，并使用引理~\ref{lem:9-1-8}，得到
\[
  \frac d{d\lambda}f^\lambda(x)
  =-\int_0^1(\partial_\lambda H_t^\lambda)(\phi_t^\lambda(x))\,dt.
\]
由刚才所证，这个表达式与 $x$ 无关。又因为每个 $H_t^\lambda$ 的均值为零，它在 $M$
上的积分为零。所以对每个 $\lambda,x$ 都有 $\partial_\lambda f^\lambda(x)=0$。
道路 $\lambda\mapsto f^\lambda(x_0)$ 的导数与
$\lambda\mapsto\mathcal A_{H^\lambda}(x^\lambda,[u^\lambda])$ 的导数相同，故引理得证。
\end{proof}

\begin{exercise}\label{ex:9-1-10}
直接计算证明，上述证明中的函数 $f^\lambda:M\to\mathbb R$ 是常值函数。

\emph{提示：}设 $\mathbb R\to M:r\mapsto x(r)$ 是任意光滑道路，并简记
$u_r:=u_{x(r)}$。利用 Cartan 公式和 Stokes 定理证明
\[
  \frac d{dr}\int_{[0,1]^2}u_r^*\omega
  =\frac d{dr}\int_0^1\left(
    H_t^0(\phi_t^0(x(r)))-H_t^\lambda(\phi_t^\lambda(x(r)))
  \right)\,dt.
\]
\end{exercise}

\begin{exercise}\label{ex:9-1-11}
设 $\mathcal P_\omega\subset\mathbb R$ 为 $\omega$ 的周期所成子群，即
\[
  \mathcal P_\omega:=\{\omega(A)\mid A\in\pi_2(M)\}.
\]
对每个 $\phi\in\operatorname{Ham}(M,\omega)$，定义一对可缩不动点
$x^0,x^1\in\operatorname{Fix}_0(\phi)$ 的\emph{\BookRecordIndex{T0347}相对辛作用量}为
\[
  \mathcal A_\phi(x^0,x^1)
  :=\int_{[0,1]^2}v^*\omega\in\mathbb R/\mathcal P_\omega,
\]
其中 $v:[0,1]\times[0,1]\to M$ 是满足
\[
  v(0,t)=x^0,\qquad v(1,t)=x^1,\qquad v(s,0)=\phi(v(s,1))
\]
的光滑映射。因此，$\mathcal A_\phi(x^0,x^1)$ 是道路 $s\mapsto v(s,0)$ 从 $x^0$
到 $x^1$ 与其在 $\phi^{-1}$ 下的像
$v(s,1)=\phi^{-1}(v(s,0))$ 之间的面积（模 $\mathcal P_\omega$）。证明
\[
  \mathcal A_\phi(x^0,x^1)
  =\mathcal A_H(x^0,[u^0])-\mathcal A_H(x^1,[u^1])
  \pmod{\mathcal P_\omega},
\]
其中 $H\in C^\infty(\mathbb R/\mathbb Z\times M)$ 是一个哈密顿函数，其\BookRecordIndex{T0158}哈密顿同痕
$t\mapsto\phi_t$ 的时间--$1$ 映射为 $\phi_1=\phi$；对 $i=0,1$，令
$x^i(t):=\phi_t(x^i)$，而 $(x^i,[u^i])$ 是 $x^i$ 到 $\widetilde{L_0M}$ 的任意提升。
由此\BookRecordIndex{T0326}推出，$\mathcal A_\phi(x^0,x^1)$ 与 $v$ 无关，并作为
$\mathbb R/\mathcal P_\omega$ 的元素有良好定义。

\emph{提示：}为了证明 $v$ 存在，任取 $L_0M$ 中从 $x^0$ 到 $x^1$ 的道路
\[
  [0,1]\longrightarrow L_0M,\qquad s\longmapsto x^s,
\]
并令 $v(s,t):=\phi_t^{-1}(x^s(t))$。再选取单位圆盘
$\mathbb D$\sourcefootnote{原书此处把圆盘记作 $B$，但本节只定义了单位圆盘
$\mathbb D$，且紧接着仍以 $e^{2\pi it}$ 参数化其边界，故统一改为 $\mathbb D$。}
上的光滑映射
\[
  [0,1]\times\mathbb D\longrightarrow M,
  \qquad (s,z)\longmapsto u^s(z),
\]
使得对所有 $s,t$ 都有 $u^s(e^{2\pi it})=x^s(t)$。利用引理~\ref{lem:9-1-8}
对曲线 $s\mapsto\mathcal A_H(x^s,[u^s])$ 求导，并证明
\[
  \frac d{ds}\mathcal A_H(x^s,[u^s])
  =-\int_0^1\omega(\partial_sv(s,t),\partial_tv(s,t))\,dt.
\]
（简记 $u(s,t):=u^s(e^{2\pi it})$，并分别对 $s,t$ 求导恒等式
$u(s,t)=\phi_t(v(s,t))$。）另一种做法是定义
$w:[0,1]^2\to M$ 为 $w(s,t):=\phi_t(v(s,1))$，证明
$u^0-u^1+v+w$ 是一个 $2$-循环（所表示的类属于 Hurewicz 同态
$\pi_2(M)\to H_2(M)$ 的像），再在这个循环上积分 $\omega$。（$w^*\omega$ 的积分
给出哈密顿项。）
\end{exercise}

\paragraph{弗洛尔方程。}
定理~\ref{thm:9-1-6} 的证明发展了定理~\ref{thm:9-1-1} 证明中的思想；不过非球面
情形没有\BookRecordIndex{T0243}冒泡，因而证明更容易。我们先把那里引入的一些概念形式化。

固定一族光滑的单参数 $\omega$-\BookRecordIndex{T0348}相容\BookRecordIndex{T0088}殆复结构
$J_t=J_{t+1}\in\mathcal J(M,\omega)$。这样的族在每个切空间
\[
  T_xL_0M=C^\infty(\mathbb R/\mathbb Z,x^*TM)
\]
上决定 $L^2$ 内积
\[
  \langle\xi,\eta\rangle
  :=\int_0^1\omega(\xi(t),J_t(x)\eta(t))\,dt.
\]
由引理~\ref{lem:9-1-8}，$\mathcal A_H$ 关于该内积的梯度为
\[
  \operatorname{grad}\mathcal A_H(x)
  =J_t(x)(\dot x-X_{H_t}(x)).
\]
因此，$\mathcal A_H$ 的一条（负）梯度流线就是满足 \emph{弗洛尔方程}
\begin{equation}\label{eq:9-1-11}
  \partial_su+J_t(u)(\partial_tu-X_{H_t}(u))=0
\end{equation}
的光滑函数 $u:\mathbb R\times\mathbb R/\mathbb Z\to M$。即使 $(M,\omega)$ 不是\BookRecordIndex{T0355}辛非
球面的，或者 $J_t$ 只被 $\omega$ \BookRecordIndex{T0372}驯服，方程~\eqref{eq:9-1-11} 仍有意义。Floer
\srcite{114} 在研究\BookRecordIndex{T0132}弗洛尔同调时引入了它。式~\eqref{eq:9-1-11} 的一个解 $u$ 的
能量定义为
\[
  E_H(u):=\frac12\int_{-\infty}^{\infty}\int_0^1
  \left(|\partial_su|_{J_t}^2
  +|\partial_tu-X_{H_t}(u)|_{J_t}^2\right)\,dt\,ds.
\]
以 $\mathcal M^{\mathrm{Floer}}(H,J)$ 表示式~\eqref{eq:9-1-11} 的可缩有限能量解之集。
下一个引理说明，若某个 $J$ 使 $\mathcal M^{\mathrm{Floer}}(H,J)\ne\varnothing$，则
$\mathcal P_0(H)\ne\varnothing$。尽管引理是对\BookRecordIndex{T0355}辛非球面流形表述的，这一结论一般也成立。

\begin{lemma}\label{lem:9-1-12}
设 $(M,\omega)$ 是紧\BookRecordIndex{T0355}辛非球面流形。设
$H:\mathbb R/\mathbb Z\times M\to\mathbb R$ 是光滑哈密顿函数，
$J_t=J_{t+1}$ 是 $M$ 上一族光滑的 $\omega$-\BookRecordIndex{T0372}驯服\BookRecordIndex{T0088}殆复结构，并且
$u\in\mathcal M^{\mathrm{Floer}}(H,J)$。则有：
\begin{enumerate}[label=\textnormal{(\roman*)}]
\item $\displaystyle\lim_{s\to\pm\infty}\sup_{t\in\mathbb R}
|\partial_su(s,t)|=0$。

\item 存在 $x,y\in\mathcal P_0(H)$，使
\[
  E_H(u)=\mathcal A_H(x)-\mathcal A_H(y),
\]
并且
\[
  \mathcal A_H(x)=\sup_{s\in\mathbb R}\mathcal A_H(u(s,\cdot)),
  \qquad
  \mathcal A_H(y)=\inf_{s\in\mathbb R}\mathcal A_H(u(s,\cdot)).
\]

\item 若 $\mathcal P_0(H)$ 是有限集，则 (ii) 中的周期解 $x,y$ 可以选成满足
\[
  \lim_{s\to-\infty}u(s,t)=x(t),
  \qquad
  \lim_{s\to\infty}u(s,t)=y(t),
\]
且收敛关于 $t$ 一致。
\end{enumerate}
\end{lemma}

\begin{proof}
先证明 $\sup_{s,t}|\partial_su(s,t)|<\infty$。因为 $E_H(u)<\infty$，存在 $T>0$，使
$u$ 在集合 $(\mathbb R\setminus[-T,T])\times[0,1]$ 上的限制之能量，小于每个全纯
球面的能量（对每个 $J_t$ 都如此）。因此，第~\ref{sec:4-2} 节中的\BookRecordIndex{T0243}冒泡论证说明函数
$(s,t)\mapsto|\partial_su(s,t)|$ 有界。

证明 (i)。反设存在序列 $s_\nu,t_\nu\in\mathbb R$ 和常数 $\delta>0$，使
$|s_\nu|\to\infty$、$0\le t_\nu\le1$，且
$|\partial_su(s_\nu,t_\nu)|\ge\delta$。由于 $|\partial_su|$ 有界，附录 B 的\BookRecordIndex{T0330}椭圆自举
分析以及 \personlater{geluomofu}{格罗莫夫}{Gromov}的图像构造（它把
式~\eqref{eq:9-1-11} 之解的图像解释为 $\widetilde J_H$-全纯曲线；见注
~\ref{rem:8-1-9}）说明，存在子列 $\nu_i$，使序列
\[
  (s,t)\longmapsto u(s_{\nu_i}+s,t)
\]
在紧集上 $C^\infty$ 收敛到式~\eqref{eq:9-1-11} 的一个解 $u^\infty(s,t)$。对每个
固定的 $T$，$u$ 在区域 $[s_{\nu_i}-T,s_{\nu_i}+T]\times[0,1]$ 上的能量趋于零，
所以极限函数 $u^\infty$ 的能量必为零。于是
$\partial_su(s_{\nu_i},t_{\nu_i})\to0$，矛盾。这证明了 (i)。

证明 (ii)。由引理~\ref{lem:9-1-8}，对 $a<b$ 有
\begin{align}
  0&\le E_H(u;[a,b]\times[0,1])\notag\\
   &=\int_a^b\int_0^1
     \omega(\partial_su,\partial_tu-X_{H_t}(u))\,dt\,ds\label{eq:9-1-12}\\
   &=\mathcal A_H(u(a,\cdot))-\mathcal A_H(u(b,\cdot)).\notag
\end{align}
只须证明，当 $s\to\pm\infty$ 时，$\mathcal A_H(u(s,\cdot))$ 的极限是
$\mathcal A_H$ 的\BookRecordIndex{T0234}临界值。为此先注意，当 $|s|\to\infty$ 时，
$|\partial_tu-X_{H_t}(u)|$ 关于 $t$ 一致趋于零。因此，对任意序列 $s_\nu\to\infty$，
都存在子列 $\nu_i$ 及 $y\in\mathcal P_0(H)$，使
$u(s_{\nu_i},t)$ 关于 $t$ 一致收敛到 $y(t)$。又因为当 $|s|\to\infty$ 时
$\partial_tu-X_{H_t}(u)\to0$，所以 $u(s_{\nu_i},\cdot)$ 在 $C^1$ 拓扑下收敛到
$y$，进而 $\mathcal A_H(u(s_{\nu_i},\cdot))\to\mathcal A_H(y)$。对
$s_\nu\to-\infty$ 的序列同样成立。这证明了 (ii)。

证明 (iii)。若 $\mathcal P_0(H)$ 有限，则式~\eqref{eq:9-1-1} 的可缩 $1$-周期解都是
孤立的。因此，极限 $x$ 与序列 $s_\nu\to-\infty$ 的选取无关，极限 $y$ 与序列
$s_\nu\to+\infty$ 的选取无关。这证明了 (iii)，也证明了引理。
\end{proof}

由引理~\ref{lem:9-1-12}，周期轨道的存在性定理可以归结为 弗洛尔方程解的存在性定理。
回忆哈密顿函数
$\mathbb R/\mathbb Z\to C^\infty(M):t\mapsto H_t$ 的 \BookRecordIndex{T0028}Hofer 范数定义为
\[
  \|H\|:=\int_0^1\left(\max_MH_t-\min_MH_t\right)\,dt.
\]
下一个定理的证明推迟到本节末尾。

\begin{theorem}\label{thm:9-1-13}
设 $(M,\omega)$ 是紧\BookRecordIndex{T0355}辛非球面流形，$H\in C^\infty(\mathbb R/\mathbb Z\times M)$，
且 $\{J_t\}_{t\in\mathbb R/\mathbb Z}$ 是 $\mathcal J_\tau(M,\omega)$ 中的光滑环路。
则对每个 $x_0\in M$，存在 $u\in\mathcal M^{\mathrm{Floer}}(H,J)$，使
$u(0,0)=x_0$、$E_H(u)\le\|H\|$，并且对每个 $s\in\mathbb R$ 都有
\begin{equation}\label{eq:9-1-13}
  -\int_0^1\max_MH_t\,dt
  \le\mathcal A_H(u(s,\cdot))
  \le-\int_0^1\min_MH_t\,dt.
\end{equation}
\end{theorem}

\begin{proof}[定理~\ref{thm:9-1-6} 的证明]
设 $t\mapsto\phi_t$ 是 $H$ 所生成的\BookRecordIndex{T0158}哈密顿同痕，并设
$\phi_1\ne\operatorname{id}$。选取 $x_0\in M$ 使 $\phi_1(x_0)\ne x_0$，再任取 $M$ 上
一族光滑的 $\omega$-\BookRecordIndex{T0372}驯服\BookRecordIndex{T0088}殆复结构
$\{J_t\}_{t\in\mathbb R/\mathbb Z}$。由定理~\ref{thm:9-1-13}，存在满足
$u(0,0)=x_0$ 和式~\eqref{eq:9-1-13} 的
$u\in\mathcal M^{\mathrm{Floer}}(H,J)$。因为 $x_0$ 不是 $\phi_1$ 的不动点，
$E_H(u)>0$。所以引理~\ref{lem:9-1-12}(ii) 中的周期解
$x,y\in\mathcal P_0(H)$ 满足定理~\ref{thm:9-1-6} 的全部要求。
\end{proof}

定理~\ref{thm:9-1-13} 的另一个推论是：紧\BookRecordIndex{T0355}辛非球面流形的哈密顿辛同胚群上的 \BookRecordIndex{T0027}Hofer
度量是非退化的。（关于 \BookRecordIndex{T0027}Hofer 度量的更多细节，见下面第~\ref{sec:9-6} 节。）以
$\phi_H:M\to M$ 表示 $H\in C^\infty(\mathbb R/\mathbb Z\times M)$ 所生成的哈密顿
辛同胚。

\begin{corollary}[\BookRecordIndex{T0027}Hofer 度量]\label{cor:9-1-14}
设 $(M,\omega)$ 是紧\BookRecordIndex{T0355}辛非球面流形，且
$\phi\ne\operatorname{id}$ 是 $M$ 的哈密顿辛同胚。则
\[
  d(\phi,\operatorname{id})
  :=\inf\bigl\{\|H\|\bigm|
  H\in C^\infty(\mathbb R/\mathbb Z\times M),\ \phi_H=\phi\bigr\}>0.
\]
\end{corollary}

\begin{proof}
选取一族光滑\BookRecordIndex{T0088}殆复结构
$\mathbb R\to\mathcal J_\tau(M,\omega):t\mapsto J_t$，使对每个 $t$ 都有
$J_{t+1}=\phi^*J_t$。考虑光滑函数 $v:\mathbb R^2\to M$ 所满足的方程
\begin{equation}\label{eq:9-1-14}
  \partial_sv+J_t(v)\partial_tv=0,
  \qquad v(s,t)=\phi(v(s,t+1)).
\end{equation}
定义 $v$ 的能量为
\[
  E(v):=\int_{-\infty}^{\infty}\int_0^1
  |\partial_sv(s,t)|_{J_t}^2\,dt\,ds.
\]
选取 $x_0\in M$ 使 $\phi(x_0)\ne x_0$。由\BookRecordIndex{T0192}紧性，
\begin{equation}\label{eq:9-1-15}
  \delta:=\inf\bigl\{E(v)\bigm|
  v\text{ 满足式~\eqref{eq:9-1-14}},\ v(0,0)=x_0,\ E(v)<\infty\bigr\}>0.
\end{equation}
现在令 $H\in C^\infty(\mathbb R/\mathbb Z\times M)$ 满足 $\phi_H=\phi$，并以
$t\mapsto\phi_t$ 表示 $H$ 所生成的\BookRecordIndex{T0158}哈密顿同痕。则 $v$ 是式~\eqref{eq:9-1-14} 的解，
当且仅当函数 $u(s,t):=\phi_t(v(s,t))$ 是式~\eqref{eq:9-1-11} 的解，其中 $J_t$
换成 $\phi_{t*}J_t$。此外，$E_H(u)=E(v)$ 且 $u(0,0)=v(0,0)$。由定理
~\ref{thm:9-1-13}，式~\eqref{eq:9-1-11} 存在有限能量解 $u$（其中 $J_t$ 换成
$\phi_{t*}J_t$），使得 $u(0,0)=x_0$ 且 $E_H(u)\le\|H\|$。于是
\[
  v(s,t):=\phi_t^{-1}(u(s,t))
\]
是式~\eqref{eq:9-1-14} 的有限能量解，且 $v(0,0)=x_0$。故
\[
  \delta\le E(v)=E_H(u)\le\|H\|.
\]
推论得证。
\end{proof}

\begin{exercise}\label{ex:9-1-15}
证明式~\eqref{eq:9-1-15} 中的 $\delta>0$。
\end{exercise}

现在来证明定理~\ref{thm:9-1-13}。与定理~\ref{thm:9-1-1} 的情形一样，我们把所需的
弗洛尔轨道构造成二维球面上受扰\BookRecordIndex{T0209}柯西--黎曼方程~\eqref{eq:9-1-4} 之解的极限。给定光滑
\BookRecordIndex{T0088}殆复结构族 $J:S^2\to\mathcal J_\tau(M,\omega)$ 及\BookRecordIndex{T0157}哈密顿扰动
$K\in\Omega^1(S^2,C_0^\infty(M))$，仍以 $\mathcal M_0(J,K)$ 表示
式~\eqref{eq:9-1-4} 的可缩解之空间。

\begin{proposition}\label{prop:9-1-16}
设 $(M,\omega)$ 是紧\BookRecordIndex{T0355}辛非球面流形，
$J\in C^\infty(S^2,\mathcal J_\tau(M,\omega))$，且
$K\in\Omega^1(S^2,C_0^\infty(M))$。则对每个 $w\in S^2$，\BookRecordIndex{T0283}求值映射
\[
  \operatorname{ev}_{w,J,K}:\mathcal M_0(J,K)\longrightarrow M,
  \qquad \operatorname{ev}_{w,J,K}(v):=v(w)
\]
都是满射。
\end{proposition}

\begin{proof}
由定理~\ref{thm:4-6-1}，对每个
$J\in\mathcal J:=C^\infty(S^2,\mathcal J_\tau(M,\omega))$ 和每个
$K\in\mathcal H:=\Omega^1(S^2,C_0^\infty(M))$，\BookRecordIndex{T0246}模空间 $\mathcal M_0(J,K)$ 都紧。
给定 $J\in\mathcal J$，以 $\mathcal H_{\mathrm{reg}}(J)$ 表示所有满足如下条件的\BookRecordIndex{T0157}哈密顿
扰动 $K$ 之集：对式~\eqref{eq:9-1-4} 的每个解 $v:S^2\to M$，竖直\BookRecordIndex{T0346}线性化算子
$D_{K,v}$ 都是满射（见第~\ref{sec:8-3} 节）。于是，当
$K\in\mathcal H_{\mathrm{reg}}(J)$ 时，\BookRecordIndex{T0246}模空间是紧定向光滑 $2n$ 维流形
（定理~\ref{thm:8-3-1}），而且它的配边类与 $J\in\mathcal J$ 及
$K\in\mathcal H_{\mathrm{reg}}(J)$ 无关（定理~\ref{thm:8-3-3}）。如命题
~\ref{prop:9-1-4} 的证明一样，当 $J=J_0$、$K=0$ 时，\BookRecordIndex{T0283}求值映射
$\operatorname{ev}_{w,J,K}:\mathcal M_0(J,K)\to M$ 的次数为一。因此，对每个
$J\in\mathcal J$ 和每个 $K\in\mathcal H_{\mathrm{reg}}(J)$，这仍成立。又因为
$\mathcal H_{\mathrm{reg}}(J)$ 在 $\mathcal H$ 中稠密（定理~\ref{thm:8-3-1}），\BookRecordIndex{T0192}紧性
蕴含：对每个 $J\in\mathcal J$、$K\in\mathcal H$ 和 $w\in S^2$，
$\operatorname{ev}_{w,J,K}$ 都是满射。
\end{proof}

\begin{proof}[定理~\ref{thm:9-1-13} 的证明]
如式~\eqref{eq:9-1-6}，令
$K_T\in\Omega^1(S^2,C_0^\infty(M))$ 为如下 $1$-形式：它到
$\mathbb R\times\mathbb R/\mathbb Z$ 的拉回是
\[
  \phi^*K_T=-\beta_T(s)H_t\,dt,
\]
其中 $\beta_T:\mathbb R\to[0,1]$ 是紧支光滑\BookRecordIndex{T0188}截断函数，在 $|s|\le T$ 时等于一，并且
在 $s\ge0$ 时非增、在 $s\le0$ 时非减。另选取光滑族
$J_T:S^2\to\mathcal J_\tau(M,\omega)$，使它在圆柱区域 $|s|\le T$ 上等于给定环路
$J_t$，并在两极邻域光滑延拓；这样的延拓因
$\mathcal J_\tau(M,\omega)$ 可缩而存在。\sourcefootnote{原书在随后直接使用了未定义的
$J_T$。这里根据中央圆柱上必须还原 弗洛尔方程所用的 $J_t$，以及\BookRecordIndex{T0372}驯服\BookRecordIndex{T0088}殆复结构空间的
可缩性，补出这一光滑延拓族。}固定 $x_0\in M$，在命题
~\ref{prop:9-1-16} 中取 $w=1\in\mathbb C\subset S^2$，得到式
~\eqref{eq:9-1-4} 的一个可缩解 $v_T:S^2\to M$，其中 $K=K_T$、$J=J_T$，且
$v_T(1)=x_0$。于是映射
\[
  u_T(s,t):=v_T(e^{2\pi(s+it)})
\]
在 $|s|\le T$ 时满足式~\eqref{eq:9-1-11}，并且 $u_T(0,0)=x_0$。此外，对
$-T\le s\le T$ 有
\begin{equation}\label{eq:9-1-16}
  -\int_0^1\max_MH_t\,dt
  \le\mathcal A_H(u_T(s,\cdot))
  \le-\int_0^1\min_MH_t\,dt.
\end{equation}
为了说明这一点，注意当 $s\ge-T$ 时，
\begin{align*}
  \frac d{ds}\mathcal A_{\beta_T(s)H}(u_T(s,\cdot))
  &=-\int_0^1|\partial_su_T(s,t)|_{J_T}^2\,dt
    -\int_0^1\dot\beta_T(s)H_t(u_T(s,t))\,dt\\
  &\le-\dot\beta_T(s)\int_0^1\max_MH_t\,dt.
\end{align*}
又注意，当 $s\to\infty$ 时，
$\mathcal A_{\beta_T(s)H}(u_T(s,\cdot))\to0$。在区间 $[s,\infty)$ 上积分，便得到
式~\eqref{eq:9-1-16} 中第一个不等式；第二个不等式同理可证。

由式~\eqref{eq:9-1-16} 和~\eqref{eq:9-1-12}，$u_T$ 在圆柱
$[-T,T]\times\mathbb R/\mathbb Z$ 上的能量不超过 $\|H\|$。由于不存在
$J$-全纯球面，我们得到无\BookRecordIndex{T0243}冒泡的\BookRecordIndex{T0192}紧性。因此，函数
$(s,t)\mapsto|\partial_su_T(s,t)|$ 在每个紧集上一致有界（见第~\ref{sec:4-2} 节）。再由定理
B.4.2（用于图像），存在序列 $T_\nu\to\infty$，使 $u_{T_\nu}$ 在每个紧集上的
$C^\infty$ 拓扑中收敛到式~\eqref{eq:9-1-11} 的一个有限能量解 $u$。极限满足
$u(0,0)=x_0$，而式~\eqref{eq:9-1-13} 由式~\eqref{eq:9-1-16} 得出。此外，每个
$u_T$ 都可缩，故 $u$ 也可缩。定理得证。
\end{proof}
